\section[Cote 46 : relever un cocycle à travers la réduction modulo $\mathfrak m$]{Cote 46 : relever un cocycle à travers la réduction modulo $\mathfrak m$\protect\verifie{Folder46}}\label{sec:46}

La cote 46, \emph{Schémas en groupes et fibrés principaux (Divers)}, quatre-vingt-quatorze pages sans pagination de sa main, réunit des liasses sur le Frobenius, les schémas de Brauer--Severi, la restriction des scalaires et, aux pages 71 à 92, des \q{Petits fourbis cohomologiques} : théorème~90 tordu, théorème de Lang, torsion des cocycles, version semi-locale. Les pages 90 à 92 en sont la dernière étape (\lot{46}{5}{90}). Soient $\mathcal O$ un anneau local, d'idéal maximal $\mathfrak m$ et de corps résiduel $k$ de caractéristique $p$, $\Gamma$ un groupe fini et $A$ une $(\mathcal O, \Gamma)$-algèbre finie, non nécessairement commutative : la marge de la page~90 cite $A = \mathrm M_n(\mathcal O)$. La page affirme que $H^1(\Gamma, A^*) \to H^1(\Gamma, (A/\mathfrak mA)^*) \times H^1(\Lambda, A^*)$ est injective, $\Lambda$ un $p$-sous-groupe de Sylow, et \q{en particulier}, quand $p$ ne divise pas l'ordre de $\Gamma$, que la réduction $H^1(\Gamma, A^*) \to H^1(\Gamma, (A/\mathfrak mA)^*)$ l'est. La lecture modernisée n'en retient que la forme \q{noyau trivial}. On démontre ici ce cas particulier, par une autre voie que la page, et montre qu'il se passe des deux hypothèses que la page discute ; puis l'injectivité elle-même, que la page affirme, et enfin la forme de Sylow, sous la forme \q{noyau trivial}.

\noindent\emph{Cocycles et cobords.} Le groupe $\Gamma$ opère sur l'anneau $A$ par automorphismes d'anneau, et donc sur le groupe $A^*$ de ses inversibles. Un $1$-cocycle est une application $c : \Gamma \to A^*$ telle que $c(gh) = c(g)\, g(c(h))$ ; c'est un homomorphisme quand l'action est triviale, et en général un homomorphisme \q{tordu} par l'action. Il est cohomologue au cocycle trivial, ou est un cobord, s'il existe $u \in A^*$ avec $c(g) = u^{-1} g(u)$ pour tout $g$. L'ensemble $H^1(\Gamma, A^*)$ des classes n'est pas un groupe quand $A^*$ n'est pas commutatif, seulement un ensemble pointé, et \q{noyau trivial} y est plus faible qu'\q{injectif}. Dire que la classe de $c$ meurt dans $H^1(\Gamma, (A/\mathfrak mA)^*)$, c'est dire que l'image $\bar c$ de $c$ modulo $\mathfrak mA$ est un cobord : $\bar c(g) = \bar v^{-1} g(\bar v)$ pour un inversible $\bar v$ de $A/\mathfrak mA$. La question est de relever ce cobord.

\noindent\emph{La voie de la page : dévisser par les puissances de $\mathfrak m$.} Supposons d'abord $\mathfrak m^N A = 0$. La page corrige le cocycle couche par couche. Si $c$ est déjà trivial modulo $\mathfrak m^n A$, on peut l'écrire $c(g) = 1 + a(g)$ avec $a(g) \in \mathfrak m^n A$. Modulo $\mathfrak m^{n+1}A$, le produit $a(g)\,g(a(h))$ disparaît, et la relation de cocycle devient additive : $a(gh) = a(g) + g(a(h))$. La couche $\mathfrak m^n A/\mathfrak m^{n+1}A$ est un groupe abélien, et même un espace vectoriel sur $k$. Sur un tel module, quand $|\Gamma|$ est inversible, tout cocycle additif est un cobord : on pose $w = -\frac{1}{|\Gamma|}\sum_{\sigma} a(\sigma)$, et la relation de cocycle donne $a(g) = g(w) - w$. Alors $1 + w$ est inversible, et $(1+w)^{-1}g(1+w) = 1 + a(g)$ modulo $\mathfrak m^{n+1}A$ : le cocycle est trivial une couche plus bas. Après $N$ couches, il est trivial. L'inversibilité de $|\Gamma|$ est là où sert \q{$p$ ne divise pas $|\Gamma|$} : dans un espace vectoriel sur $k$, diviser par $|\Gamma|$ demande que $|\Gamma| \neq 0$ dans $k$. Pour le cas général, $A$ fini sur $\mathcal O$ sans $\mathfrak m$ nilpotent, la lecture passe à la limite par Artin--Rees.

\noindent\emph{La voie suivie ici : une seule moyenne.} La démonstration du théorème~90 de Hilbert fournit un raccourci. Pour un cocycle $c$ et un élément quelconque $w$ de $A$, posons $b = \sum_{\sigma \in \Gamma} c(\sigma)\,\sigma(w)$. Alors
\[
c(g)\, g(b) = \sum_\sigma c(g)\, g(c(\sigma))\, g\sigma(w) = \sum_\sigma c(g\sigma)\, (g\sigma)(w) = b,
\]
puisque $\sigma \mapsto g\sigma$ permute $\Gamma$. Si $b$ est inversible, cela s'écrit $c(g) = b\, g(b)^{-1}$, et $c$ est un cobord, avec $u = b^{-1}$. Tout revient à choisir $w$ pour que $b$ soit inversible. Si $c$ est trivial modulo $\mathfrak mA$, on prend $w = 1$ : $b = \sum_\sigma c(\sigma)$ est congru à $|\Gamma|$, donc s'écrit $|\Gamma|(1 + x)$ avec $x \in \mathfrak mA$, et il est inversible dès que $|\Gamma|$ l'est et que les éléments $1 + x$, $x \in \mathfrak mA$, le sont. Dans le cas général, on relève $\bar v^{-1}$ en un $w \in A$ ; alors $b \equiv \sum_\sigma \bar v^{-1}\sigma(\bar v)\,\sigma(\bar v^{-1}) = |\Gamma|\, \bar v^{-1}$, et $b = |\Gamma| w (1 + x)$ avec $x \in \mathfrak mA$. Les couches de la page sont ainsi traitées d'un coup : la moyenne se fait dans $A^*$ et non dans chaque quotient abélien.

Reste à voir que $1 + \mathfrak mA$ est fait d'inversibles, et que $w$ l'est. Quand $\mathfrak m^NA = 0$, un élément $x$ de $\mathfrak mA$ vérifie $x^N = 0$, et $1 + x$ a pour inverse $1 - x + x^2 - \cdots \pm x^{N-1}$. Quand $A$ est fini sur $\mathcal O$, c'est le lemme de Nakayama : la multiplication à gauche par $1 + x$ est une application $\mathcal O$-linéaire de $A$ dans lui-même, surjective modulo $\mathfrak mA$ puisqu'elle y est l'identité, donc surjective ; $1 + x$ a un inverse à droite, de même un inverse à gauche, et il est inversible. Un élément $w$ inversible modulo $\mathfrak mA$, avec $ww' = 1 + x$ et $w'w = 1 + y$, $x, y \in \mathfrak mA$, est alors inversible des deux côtés.

\begin{encadre}
\begin{enonce}[D'après les pages 90 à 92, forme reconstituée par l'édition, cas où $p$ ne divise pas l'ordre de $\Gamma$\piste{46-sylow-reduction-criterion}]
Soient $\mathcal O$ un anneau local commutatif, de corps résiduel de caractéristique $p$, $\Gamma$ un groupe fini d'ordre premier à $p$, et $A$ une $\mathcal O$-algèbre, non nécessairement commutative, sur laquelle $\Gamma$ opère par automorphismes d'anneau. On suppose $\mathfrak m^N A = 0$ pour un entier $N$, ou $A$ fini sur $\mathcal O$. Alors le noyau de $H^1(\Gamma, A^*) \to H^1(\Gamma, (A/\mathfrak mA)^*)$ est trivial : un cocycle dont la réduction modulo $\mathfrak mA$ est un cobord est un cobord. Plus généralement, il en est ainsi de $H^1(\Gamma, A^*) \to H^1(\Gamma, B^*)$ pour tout morphisme d'anneaux surjectif $A \to B$ compatible à $\Gamma$ et de noyau contenu dans $\mathfrak mA$.
\end{enonce}
\end{encadre}

\begin{proof}
Soit $c$ un cocycle et $\bar v$ un inversible de $B$ avec $\bar c(g) = \bar v^{-1} g(\bar v)$. Comme $|\Gamma|$ n'est pas divisible par $p$, son image dans le corps résiduel n'est pas nulle, et $|\Gamma|$ est inversible dans $\mathcal O$, donc dans $A$. Les éléments $1 + x$, $x \in \mathfrak mA$, sont inversibles, par nilpotence ou par Nakayama. Relevons $\bar v^{-1}$ et $\bar v$ en $w$ et $w'$ ; $ww'$ et $w'w$ sont de la forme $1 + x$, donc $w$ est inversible. L'élément $b = \sum_\sigma c(\sigma)\,\sigma(w)$ a pour image $\sum_\sigma \bar v^{-1}\,\sigma(\bar v)\,\sigma(\bar v^{-1}) = |\Gamma|\,\bar v^{-1}$, la même que $|\Gamma|\, w$. Donc $b = m(1 + m^{-1}x)$, avec $m = |\Gamma|\,w$ inversible et $x$ dans le noyau, et $b$ est inversible. L'identité $c(g)\,g(b) = b$ donne $c(g) = u^{-1} g(u)$ avec $u = b^{-1}$.
\end{proof}

\begin{contrexemple*}[De l'édition : l'hypothèse sur $p$ est nécessaire]
Soit $\Gamma = \mathbb Z/2$, opérant trivialement sur $A = \mathcal O = \mathbb Z/4$, où $\mathfrak m = (2)$, $\mathfrak m^2 = 0$ et $p = 2 = |\Gamma|$. L'application $c$ qui envoie l'élément non trivial de $\Gamma$ sur $-1$ est un cocycle, trivial modulo $\mathfrak m$, puisque $-1 \equiv 1 \pmod 2$. Ce n'est pas un cobord : l'action étant triviale et $A$ commutatif, $u^{-1}g(u) = 1$ pour tout $u$, et $-1 \neq 1$ dans $\mathbb Z/4$.
\end{contrexemple*}

\begin{remarque}
La démonstration n'utilise que deux choses : $|\Gamma|$ inversible dans $A$, et $1 + x$ inversible pour $x$ dans le noyau de $A \to B$. Elle ne se sert ni de l'hypothèse du texte, que $\Gamma$ opère trivialement sur $A/\mathfrak mA$, ni de celle de la marge, que $A^\Gamma \to A/\mathfrak mA$ est surjectif, ni d'une action de $\Gamma$ sur $\mathcal O$, ni de la noethérianité, et le cas fini se passe d'Artin--Rees. La page en a besoin parce qu'elle vise la forme de Sylow : la restriction $H^1(\Gamma, A) \to H^1(\Lambda, A)$, injective puisque $[\Gamma : \Lambda]$ est inversible, fournit un $w \in A$ quelconque avec $a(g) = g(w) - w$, et la surjectivité de $A^\Gamma \to A/\mathfrak mA$ sert à corriger ce $w$ par un invariant pour que $\varepsilon(w) = 1$, donc pour que $w$ soit inversible. Quand $\Lambda$ est trivial, la moyenne donne directement un $w$ dans $\mathfrak m^nA$, et la correction n'a pas lieu d'être.
\end{remarque}

\begin{remarque}
La marge qui porte la surjectivité de $A^\Gamma \to A/\mathfrak mA$ est écrite en huit lignes obliques dans la marge gauche de la page~90, reliée par un double trait à \q{Supposons de plus que} ; la transcription y relève une douzaine de mots illisibles et lit \q{plus généralement supposons} avec doute, \q{surjectif} étant souligné deux fois. L'hypothèse dont la forme de Sylow dépend repose donc sur une lecture fragile, que seul le fac-similé peut confirmer. Le cas démontré ici n'en dépend pas.
\end{remarque}

\begin{remarque}
Pour $A$ commutatif, l'énoncé découle du lemme du §\,3.6 de Lorenz \cite{lorenz1999}, lu dans sa prépublication (arXiv:math/9808032) ; la version publiée n'a pas été consultée, et sa numérotation peut différer. Lorenz suppose $S$ commutatif et la trace $S \to S^G$ surjective, ce qu'entraîne $|G|$ inversible, et montre que $H^1(G, \mathrm{GL}_n(S)) \to H^1(G, \mathrm{GL}_n(S/\mathrm{rad}\, S))$ est de noyau trivial, par les modules projectifs et Nakayama. Comme $\mathfrak mA$ est contenu dans le radical de Jacobson de $A$ quand $A$ est fini sur $\mathcal O$, le cas $n = 1$ donne le cas commutatif de l'énoncé encadré. Le cas d'une algèbre non commutative à action quelconque, celui de la marge ($A = \mathrm M_n(\mathcal O)$ avec une action qui n'est pas forcément coefficient par coefficient), n'y est pas sous cette forme ; l'argument de moyenne est classique, et l'énoncé n'est pas présenté comme nouveau.
\end{remarque}

\noindent\emph{De noyau trivial à injectif.} Dans un ensemble pointé, un noyau trivial ne dit pas que deux classes de même image sont égales. Les pages~86 à~89 passent de l'un à l'autre en \emph{tordant} l'action par un cocycle $c'$ : on fait opérer $\Gamma$ sur $A$ par $\gamma \star a = c'(\gamma)\,\gamma(a)\,c'(\gamma)^{-1}$, et dire que $c$ et $c'$ sont cohomologues revient à dire qu'un cocycle de l'action tordue est un cobord. On peut éviter de changer d'action en faisant la moyenne directement avec les deux cocycles.

\begin{encadre}
\begin{enonce}[Grothendieck, page 90, injectivité, cas où $p$ ne divise pas l'ordre de $\Gamma$\piste{46-semilocal-hilbert90-inertia}]
Sous les hypothèses de l'énoncé précédent, avec $A$ fini sur $\mathcal O$, l'application $H^1(\Gamma, A^*) \to H^1(\Gamma, B^*)$ est injective : si deux cocycles $c, c' : \Gamma \to A^*$ ont des réductions cohomologues dans $B^*$, ils sont cohomologues, c'est-à-dire que $c(g) = u^{-1}\,c'(g)\,g(u)$ pour un $u \in A^*$.
\end{enonce}
\end{encadre}

\begin{proof}
Pour $w \in A$, posons $u = \sum_\sigma c'(\sigma)\,\sigma(w)\,c(\sigma)^{-1}$. La relation de cocycle donne $g(c(\sigma))^{-1} = c(g\sigma)^{-1} c(g)$, d'où
\[
c'(g)\,g(u) = \sum_\sigma c'(g)\,g(c'(\sigma))\,(g\sigma)(w)\,g(c(\sigma))^{-1} = \sum_\sigma c'(g\sigma)\,(g\sigma)(w)\,c(g\sigma)^{-1}\,c(g) = u\,c(g).
\]
Si $u$ est inversible, cela s'écrit $c(g) = u^{-1} c'(g)\,g(u)$. Soit $v \in B^*$ avec $\bar c(g) = v^{-1}\,\bar c'(g)\,g(v)$ ; relevons $v$ en $w$, inversible comme plus haut. Alors $\bar c'(\sigma)\,\sigma(v)\,\bar c(\sigma)^{-1} = v$ pour tout $\sigma$, de sorte que $u$ a la même image que $|\Gamma|\,w$, et $u = |\Gamma| w (1 + x)$ avec $x$ dans le noyau : $u$ est inversible.
\end{proof}

\noindent Avec $c' = 1$, on retrouve la moyenne de l'énoncé précédent. La moyenne $u$ est celle du théorème~90 pour l'action tordue par $c'$, écrite sans changer d'action.

\noindent\emph{La forme de Sylow.} Reste le facteur $H^1(\Lambda, A^*)$ de la page~90. Quand $p$ divise l'ordre de $\Gamma$, la moyenne sur $\Gamma$ n'est plus possible, mais la moyenne sur les classes à gauche $g\Lambda$ l'est encore, puisque leur nombre $[\Gamma : \Lambda]$ est premier à $p$. Il faut pour cela que le terme $c(g)\,g(z)$ ne dépende que de la classe $g\Lambda$, c'est-à-dire que $c(\lambda)\,\lambda(z) = z$ pour $\lambda \in \Lambda$ : la relation de cocycle donne alors $c(g\lambda)\,(g\lambda)(z) = c(g)\,g\bigl(c(\lambda)\,\lambda(z)\bigr) = c(g)\,g(z)$. C'est là que servent les deux hypothèses de la page. Que $c$ soit un cobord sur $\Lambda$, $c(\lambda) = u^{-1}\lambda(u)$, permet de prendre $z = u^{-1}w$ avec $w$ invariant par $\Lambda$ ; que $w$ puisse être choisi de réduction prescrite modulo $\mathfrak mA$, c'est la surjectivité de la marge.

\begin{encadre}
\begin{enonce}[D'après la page 90 et sa marge, forme \q{noyau trivial}\piste{46-sylow-reduction-criterion}]
Soient $\mathcal O$ un anneau local commutatif, de corps résiduel de caractéristique $p$, $\Gamma$ un groupe fini, $\Lambda$ un sous-groupe d'indice premier à $p$, par exemple un $p$-sous-groupe de Sylow, et $A$ une $\mathcal O$-algèbre finie, non nécessairement commutative, sur laquelle $\Gamma$ opère par automorphismes d'anneau. On suppose que $A^\Lambda \to A/\mathfrak mA$ est surjectif, ce qu'entraîne la surjectivité de $A^\Gamma \to A/\mathfrak mA$ demandée par la marge. Alors une classe de $H^1(\Gamma, A^*)$ triviale dans $H^1(\Gamma, (A/\mathfrak mA)^*)$ et dans $H^1(\Lambda, A^*)$ est triviale. Il en est de même pour tout morphisme d'anneaux surjectif $\pi : A \to B$ compatible à $\Gamma$, de noyau contenu dans $\mathfrak mA$, tel que $A^\Lambda \to B$ soit surjectif.
\end{enonce}
\end{encadre}

\begin{proof}
Soient $u \in A^*$ avec $c(\lambda) = u^{-1}\lambda(u)$ pour $\lambda \in \Lambda$, et $\bar v \in B^*$ avec $\bar c(g) = \bar v^{-1} g(\bar v)$. Choisissons $w \in A^\Lambda$ de réduction $\pi(u)\,\bar v^{-1}$, et posons $z = u^{-1}w$. Pour $\lambda \in \Lambda$,
\[
c(\lambda)\,\lambda(z) = u^{-1}\lambda(u)\,\lambda(u^{-1})\,\lambda(w) = u^{-1} w = z ,
\]
de sorte que $b = \sum_{g\Lambda \in \Gamma/\Lambda} c(g)\,g(z)$ est bien défini. Comme plus haut, $c(\gamma)\,\gamma(b)$ est la même somme, réindexée par $g\Lambda \mapsto \gamma g\Lambda$, d'où $c(\gamma)\,\gamma(b) = b$. La réduction de $z$ est $\bar v^{-1}$, celle de chaque terme est $\bar v^{-1}g(\bar v)\,g(\bar v^{-1}) = \bar v^{-1}$, et celle de $b$ est $[\Gamma : \Lambda]\,\bar v^{-1}$. L'indice est inversible dans $A$, $z$ l'est puisque sa réduction l'est, et $b = m(1 + m^{-1}x)$ avec $m = [\Gamma : \Lambda]\,z$ et $x$ dans le noyau : $b$ est inversible, et $c(g) = b\,g(b)^{-1}$.
\end{proof}

Pour $\Lambda = \Gamma$, l'hypothèse sur $\Lambda$ dit déjà que $c$ est un cobord, et l'énoncé est vide ; pour $\Lambda$ trivial, $A^\Lambda = A$ se surjecte toujours sur $A/\mathfrak mA$, et l'on retrouve le premier énoncé. La démonstration n'utilise pas l'hypothèse du texte que $\Gamma$ opère trivialement sur le corps résiduel. Elle donne le noyau trivial et non l'injectivité : pour comparer deux cocycles $c$ et $c'$ égaux sur $\Lambda$, il faudrait un $w$ qui commute aux $c(\lambda)$, ce que l'invariance par $\Lambda$ ne fournit pas quand $A$ n'est pas commutatif. Pour $A$ commutatif, $H^1(\Gamma, A^*)$ est un groupe et le noyau trivial suffit.

\begin{sloppypar}
\noindent\emph{Portée de la vérification.} Lean définit les cocycles et les cobords à valeurs dans $A^*$ pour un anneau quelconque muni d'une action par automorphismes (\texttt{IsCocycle}, \texttt{IsCoboundary}), sans recourir à la cohomologie des groupes de mathlib, qui est abélienne. Il démontre l'identité de moyenne (\texttt{mul\_smul\_sum}), sa conséquence (\texttt{isCoboundary\_of\_isUnit\_sum}), le noyau trivial pour un morphisme surjectif compatible à $\Gamma$ quand $1 + \mathrm{Ker}$ est fait d'inversibles (\texttt{isCoboundary\_of\_map\_isCoboundary}), la nilpotence de $\mathfrak mA$ quand $\mathfrak m^NA = 0$ et l'inversibilité de $1 + \mathfrak mA$ pour $A$ fini, par Nakayama (\texttt{isNilpotent\_of\_mem\_smul\_top}, \texttt{isUnit\_one\_add\_of\_mem\_smul\_top}), puis les deux formes de l'énoncé encadré (\texttt{ker\_trivial\_of\_pow\_smul\_eq\_bot}, \texttt{ker\_trivial}), l'injectivité (\texttt{mul\_smul\_sum\_twist}, \texttt{isCohomologous\_of\_map}, \texttt{injective}) et le contre-exemple (\texttt{not\_isCoboundary\_signCocycle}). La forme de Sylow, ajoutée après l'étiquette dans \href{\leanblobmain/Folder46.lean}{\texttt{Folder46.lean}}, y est démontrée par la moyenne sur les classes (\texttt{mul\_smul\_sum\_cosets}, \texttt{isCoboundary\_of\_restrict}), puis sur un anneau local pour un sous-groupe d'indice premier à $p$ et pour un sous-groupe de Sylow (\texttt{ker\_trivial\_restrict}, \texttt{ker\_trivial\_sylow}). L'hypothèse \q{$p$ ne divise pas $|\Gamma|$} y est écrite avec la caractéristique du corps résiduel de mathlib. Ne sont pas formalisés l'injectivité dans la forme de Sylow pour $A$ non commutatif, la base semi-locale des pages~84 à~89 et le passage à la limite par Artin--Rees, que le cas fini rend inutile.
\end{sloppypar}

N'est pas démontré ici, de la même cote :

\begin{encadre}
\begin{theoreme*}[Grothendieck, pages 84 à 89, Hilbert~90 semi-local et inertie ; voir aussi Lorenz 1999\piste{46-semilocal-hilbert90-inertia}]\nonverifie
Soient $\mathcal O$ un anneau semi-local noethérien, $\Gamma$ un groupe fini qui opère fidèlement et transitivement sur ses idéaux maximaux, et $A$ une $\mathcal O$-algèbre finie, non nécessairement commutative, munie d'une action semi-linéaire de~$\Gamma$. Les applications $H^1(\Gamma, A^*) \to H^1(\Gamma_d, A_{\mathfrak m}^*) \to H^1(\Gamma_i, A_{\mathfrak m}^*)$ vers les groupes de décomposition et d'inertie sont injectives comme applications d'ensembles pointés, et non seulement de noyau trivial.
\end{theoreme*}
\end{encadre}
{\small La finitude de $A$ sur $\mathcal O$ est ajoutée par l'édition, et un pas que la page admet (\q{on sait que}) peut demander une hypothèse de normalité. Lorenz considère ces applications pour un anneau commutatif et n'en donne que le noyau trivial. Le cas d'un anneau local et d'un groupe d'ordre premier à $p$, où l'injectivité est démontrée plus haut, n'en est qu'une partie : la base semi-locale, la complétion et les groupes de décomposition et d'inertie n'y figurent pas. Mathlib n'a pas la cohomologie non abélienne $H^1(\Gamma, A^*)$. On se reportera aux notes de Grothendieck (\lot{46}{5}{84}).\par}

